\documentclass[12pt,a4paper]{article}
\usepackage[utf8]{inputenc}
\usepackage[T1]{fontenc}
\usepackage{times}
\usepackage[margin=2.5cm]{geometry}
\usepackage{float}
\usepackage{booktabs}
\usepackage{makecell}
\usepackage{caption}
\usepackage{graphicx}
\usepackage[british]{babel}
\usepackage{csquotes}
\usepackage[style=apa,backend=biber,natbib=true]{biblatex}
\addbibresource{references.bib}
\begin{document}

\noindent\textbf{Supplementary material}

\bigskip
\noindent{\large\bfseries Stagnation before exit: performance progression and retention in youth track and field\par}

\medskip
\noindent Anonymised manuscript. Table and section numbers follow the order of first citation in the main text.

\begin{table}[H]
\centering
\caption*{\textbf{Table S1.} Sample flow.}
\small
\begin{tabular}{lr}
\toprule
Step & n \\
\midrule
Athletes born 1998--2006 in the registry extract & 32\,447 \\
With $\geq$1 approved result at age 13--14 & 12\,938 \\
Also competing at ages 15--16 (landmark condition) & 4\,910 \\
With known sex & 4\,909 \\
\quad born 2006 (used only for the activity-at-$\geq$18 outcome) & 434 \\
Born 1998--2005 & 4\,475 \\
\quad excluded: no scorable main-group points in both windows & 578 \\
\quad excluded: sex $\times$ event-group cells with $<$20 athletes & 20 \\
Primary sample (born 1998--2005) & 3\,877 \\
\quad active seniors & 654 \\
\quad with exact birth date (H2 analyses) & 3\,850 \\
Exit-aligned trajectory subsamples (performance / volume) & 1\,116 / 1\,523 \\
Inclusive best-group exposure sample & 3\,982 \\
Activity-at-$\geq$18 outcome sample (primary sample plus born 2006) & 4\,244 \\
\bottomrule
\end{tabular}

\vspace{0.2em}
\parbox{0.8\textwidth}{\footnotesize The completed STROBE checklist accompanies the submission; the time-stamped analysis plan (committed before estimation) is available with the analysis code. Abbreviation: STROBE, Strengthening the Reporting of Observational Studies in Epidemiology.}
\end{table}

\begin{table}[H]
\centering
\caption*{\textbf{Table S2.} Robustness of the progression--retention association. OR (95\,\% CI) per SD of the exposure named (per 100 points for the earlier level given the later level); outcome is senior activity unless stated. Abbreviations: CI, confidence interval; COVID-19, coronavirus disease 2019; GEE, generalised estimating equations; OR, odds ratio; SD, standard deviation; SE, standard error.}

\small
\begin{tabular}{lccc}
\toprule
Analysis & n & Events & OR (95\,\% CI) \\
\midrule
Primary: $\Delta$ mean of 3 best, main event group & 3\,808 & 648 & 2.02 (1.73--2.37) \\
Exposure: season best (maximum) & 3\,808 & 648 & 1.93 (1.66--2.25) \\
Exposure: mean of all results & 3\,808 & 648 & 1.59 (1.39--1.82) \\
Exposure: wind-legal marks only & 3\,761 & 646 & 2.00 (1.71--2.34) \\
Exposure: best event group at 15--16, inclusive sample & 4\,071 & 663 & 1.83 (1.61--2.07) \\
Specification check: $\Delta$ within ages 13$\to$14 & 2\,735 & 561 & 1.00 (0.88--1.13) \\
Excluding cohorts born 2004--2005 (COVID window) & 2\,894 & 498 & 1.95 (1.63--2.33) \\
Progression $\times$ COVID-window interaction & 3\,808 & 648 & 1.22 (0.91--1.63) \\
$\geq$2 scored results in both windows & 2\,997 & 625 & 2.18 (1.84--2.60) \\
Outcome: $\geq$1 result at age $\geq$20 & 3\,808 & 749 & 1.82 (1.58--2.09) \\
Outcome: $\geq$2 results in a year, ages 20--21 & 3\,808 & 612 & 2.07 (1.76--2.44) \\
Outcome: active at $\geq$18 (cohorts incl.\ 2006) & 4\,172 & 1\,640 & 1.35 (1.22--1.50) \\
Cohort A (born 1998--2001) & 1\,829 & 320 & 2.00 (1.60--2.51) \\
Cohort B (born 2002--2005) & 1\,979 & 328 & 2.03 (1.62--2.54) \\
Club-clustered robust SE & 3\,808 & 648 & 2.02 (1.70--2.41) \\
GEE, exchangeable club correlation & 3\,808 & 648 & 2.01 (1.70--2.38) \\
E-value for the primary OR (CI bound) &  &  & 2.84 (2.45) \\
\bottomrule
\end{tabular}

\end{table}


\section*{S3. Deviations from the analysis plan}

\noindent the H3 interval is a DeLong interval rather than a paired bootstrap; club-grouped cross-validation folds, the exposure defined as residuals against panel growth curves, and a random club intercept were not run (club-clustered standard errors and GEE take the place of the last); the Cox model omits cohort and club size; the cohort replication was run for the H1 odds ratio only, not for H2, H3 or the AUC gain; region, planned as a covariate, is not recorded (see Covariates); and the switcher analysis takes baseline level and result counts from the best event group, where the plan took them from the main group (the planned version is shown alongside in Table S2).

\printbibliography[heading=none]

\end{document}
